\documentclass[14pt,a4paper]{extarticle}
\usepackage[margin=20mm]{geometry}
\usepackage{fontspec}
\setmainfont{Times New Roman}
\setsansfont{Arial}
\setmonofont{Consolas}[Scale=MatchLowercase]
\usepackage{polyglossia}
\setdefaultlanguage{russian}
\setotherlanguage{english}
\usepackage{amsmath,booktabs,array}
\usepackage[hyphens]{url}
\usepackage[hidelinks]{hyperref}
\hypersetup{pdftitle={Инструментальная маршрутизация языковых моделей в аграрных задачах},pdfauthor={Антоний Сарычев; Денис Шестернин},pdfsubject={Рабочий исследовательский черновик; человеческая проверка не завершена}}
\setlength{\parindent}{1.25cm}
\setlength{\parskip}{0pt}
\setlength{\emergencystretch}{2em}
\sloppy
\begin{document}
\begin{center}\bfseries
Инструментальная маршрутизация локальных и облачных языковых моделей в аграрных задачах
\end{center}

Антоний Сарычев, Денис Шестернин


РГАУ-МСХА имени К.А. Тимирязева, Москва, Россия


Рабочий исследовательский черновик от 3 октября 2026 года. Не готов к подаче: человеческая проверка эталонов, согласование авторами и новая независимая оценка не завершены. Лицензия статьи не выбрана. Имена и порядок взяты из прямого запроса авторов; отчества, email, ORCID и руководитель не добавлены.


УДК 004.8 --- рабочий класс «Искусственный интеллект», TO BE CONFIRMED BY VENUE.


\section*{Аннотация}

На 150 русских и английских аграрных задачах исследована маршрутизация между расчётными инструментами и языковыми моделями. Локальный каскад сократил облачные вызовы со 100 до 20, но не подтвердил сохранение качества. Выявлены ошибки проверки составных ответов.


Ключевые слова: языковые модели; маршрутизация; агрометеорология; детерминированные инструменты; локальный вывод; оценка качества.


\section*{1. Введение}

Аграрная система поддержки решений получает неоднородные запросы. Объяснение по справочному документу, расчёт накопленного тепла, интерпретация заданного порога и запрос с отсутствующим измерением требуют разных способов обработки. Генерация связного текста не гарантирует правильного вычисления. При неизвестной влажности или неполном описании поля обоснованным результатом может быть запрос данных, а при отсутствии подтверждающего материала --- отказ от содержательного ответа.


Использование локальной модели на потребительском компьютере потенциально уменьшает обращения к внешнему сервису. Однако снижение этих обращений может происходить уже за счёт обычного численного инструмента. Поэтому сравнение «локальная модель против облачной» недостаточно: нужно отдельно измерить вклад детерминированных расчётов и добавочный эффект локальной генерации после их применения. Кроме качества следует учитывать задержку, ошибки сервиса и фактическое число всех внешних запросов.


Цель работы --- проверить этот добавочный эффект на контролируемой смеси аграрных и агрометеорологических задач. Основное сравнение сопоставляет инструменты с облачной моделью и тот же набор инструментов с каскадом локальной и облачной моделей. Работа представляет эмпирическое исследование конкретной конфигурации. Разработка нового общего алгоритма маршрутизации и доказательство пригодности для производственных рекомендаций не заявляются.


\section*{2. Связанные исследования}

FrugalGPT и AutoMix изучают экономичные каскады моделей и выборочную эскалацию [1, 2]. Hybrid LLM и RouteLLM рассматривают маршрутизацию по запросу, включая оценку ожидаемого качества и обучение по предпочтениям [3, 4]. Объединённая постановка маршрутизации и каскадирования показывает необходимость различать решение до генерации и решение после получения ответа [5]. RouterBench предлагает сравнительную оценку нескольких моделей и компромиссов между затратами и результативностью [6]. Следовательно, сам принцип выбора модели не является новизной настоящей работы.


Оценка неопределённости и отказа от ответа также имеет самостоятельную литературу. Semantic Entropy Probes исследует обнаружение галлюцинаций [7], а обзор abstention систематизирует условия отказа и способы его оценки [8]. Эти методы не проверяют автоматически выполнение каждого условия аграрной задачи. Наш эксперимент использовал более узкую проверку буквальной опоры в контексте; её ограничения анализируются по сохранённым ответам.


В сельском хозяйстве PLLaMa, AgriGPT и AgroLLM исследуют специализированные модели и перенос знаний [9--11]. AgriAgent и AgroTools рассматривают планирование и работу с инструментами [12, 13]. Общие benchmarks tau-bench и ToolACE показывают значимость корректного взаимодействия с инструментом и соблюдения структурированного контракта [14, 15]. Работа о BitNet характеризует иной способ уменьшения ресурсов модели [16]; она не доказывает качество нашего GPU-каскада. Gemma 4 Technical Report описывает семейство выбранной локальной модели [17], а AgXQA предоставляет предметный ориентир для agricultural QA [18].


Предполагаемый вклад состоит в разделении инструментального и модельного эффектов, двуязычной контролируемой оценке, измерении запросов и задержек на домашнем GPU и сохранении отрицательного результата. Это ограниченная предметная проверка существующих архитектур. Научная библиография содержит 18 работ: для 17 сохранено чтение релевантных разделов, для AgXQA проверены метаданные; полного прочтения всех статей не заявляем.


\section*{3. Методика}

Сравниваются пять политик. LOCAL передаёт запрос локальной модели; CLOUD --- облачной. TOOLS+CLOUD сначала проверяет доступность детерминированного расчёта и полноту обязательных входов, затем вызывает облако. Hybrid TOOLS\ensuremath{\to}LOCAL\ensuremath{\to}CLOUD использует тот же инструментальный этап, затем локальный ответ и проверку опоры; при непринятом ответе происходит облачная генерация. OFFLINE сохраняет инструменты и локальную стадию, но при недостаточной опоре завершает обработку без облака. Ответ задаётся JSON-полями action, answer, numeric, unit и tool.


Инструменты вычисляют сумму эффективных температур GDD, дефицит давления водяного пара VPD, референсную эвапотранспирацию ET0 и упрощённый дефицит осадков. GDD определяется как сумма положительных разностей между средней суточной температурой и заданной базой, без верхнего отсечения. VPD оценивается по одной температуре и относительной влажности. ET0 реализует уравнения FAO-56 с явно переданными радиацией, ветром и давлениями. Дефицит осадков равен положительной части разности заданной цели и суммы осадков. Последняя цель условна и не является рекомендацией нормы полива.


Инструмент не подставляет отсутствующие измерения. Недостающие входы, несогласованные температурные ряды, недопустимые значения или единицы приводят к запросу уточнения. Возможность инструмента определяется публичными входами, без передачи gold или категории оценивания маршрутизатору. Контролируемые пороговые и составные задания в исторической политике остаются модельными: исходный gate не является универсальным исполнителем правил.


Выборка HOME150 включает 40 QA, 30 численных задач, по 20 случаев неполных данных, выбора инструмента, интерпретации правил и составного рассуждения. В каждом языке 75 строк. Зафиксированы 25 групп документов или сценариев, однако 110 строк (73,3\%) синтетические. QA опирается на пять документов. Переводные пары и варианты одного документа зависимы; число строк нельзя приравнивать к числу независимых источников. Это не репрезентативная выборка запросов фермеров Центральной России.


Правила оценщика зафиксированы до вывода и применяются после получения ответа: сравниваются action, точный span для QA, метка для правил и название инструмента. Численный допуск реализует math.isclose: абсолютная разность не больше максимума из 0,02 и 1\% от максимальной абсолютной величины ответа и эталона, с проверкой единицы. Ошибки транспорта, структуры и незавершённые ответы остаются в знаменателе. Дополнительный token F1 описывает QA, но не заменяет основной exact score. Прокси критической ошибки охватывает только неверный ANSWER в численной задаче или задаче неполных данных; он не является полной оценкой агрономического риска.


\section*{4. Условия и протокол эксперимента}

Эксперимент выполнен на AMD Ryzen 7 8700F (8 ядер, 16 потоков), 31,62 GiB RAM и NVIDIA RTX 5060 Ti с 16 311 MiB VRAM. Локальная Gemma \texttt{gemma-4-26b-a4b-it}, IQ4\_NL\_XL, обслуживается LM Studio; контекст 4096, параллелизм один, максимум 512 выходных токенов, temperature 0 и seed 42. Запрошенная доля GPU offload --- 0,70; фактическое количество выгруженных слоёв неизвестно. Размер основного GGUF --- 14 556 687 680 байт. Активное число параметров MoE нельзя приравнивать к объёму памяти весов.


Облачная модель --- Mercury 2.5 через Inception, instant reasoning, лимит 512 токенов и JSON-object. Во всех 258 возвращённых ответах temperature 0 была заменена провайдером на 1. Для 12 таймаутов нет ответа и сведений об эффективной температуре. Поэтому детерминированность облачной генерации не заявляется. Семена и параметры локальной модели также не гарантируют побитового воспроизведения на другом оборудовании.


Фактический freeze состоялся 2 октября 2026 в 04:34:45 UTC. Дата 1 октября в старом имени файла не означает более раннего freeze. Выполнены 750 отдельных решений: 150 задач \ensuremath{\times} 5 политик, с перемешанным чередованием политик. Сбой замены manifest в Windows прервал первую часть после шести задач; продолжение выполнило оставшиеся 144 без изменения входов, модели, промптов и оценщика. Уже записанные ответы не повторялись. Обе части и сведения о прерывании сохранены.


Основной контраст --- качество hybrid минус TOOLS+CLOUD. Успех требует одновременно: нижняя 95\%-я граница разности не ниже -5 процентных пунктов; сокращение облачных вызовов не меньше 10 пунктов от всех запросов; верхняя граница разности прокси критических ошибок не выше нуля; p95 задержки hybrid не больше 10 секунд. Это исследовательские пороги, а не агрономические стандарты. Применён парный кластерный bootstrap: 5000 повторов по 25 зафиксированным группам. Зависимость более высокого уровня, общих документов и шаблонов может оставаться.


\section*{5. Результаты замороженного исследования}

\begin{center}\fontsize{12}{14}\selectfont
\begin{tabular}{lrrrr}\toprule
Политика & Верно /150 & Облако & Средняя, с & p95, с \\
\midrule
LOCAL & 102 & 0 & 1,508 & 2,349 \\
CLOUD & 130 & 150 & 4,758 & 38,115 \\
TOOLS+CLOUD & 139 & 100 & 2,682 & 3,340 \\
Hybrid & 133 & 20 & 0,844 & 1,561 \\
OFFLINE & 114 & 0 & 0,655 & 1,458 \\
\bottomrule\end{tabular}
\end{center}

Hybrid сокращает генеративные облачные вызовы на 80\% относительно TOOLS+CLOUD (100\ensuremath{\to}20), но качество уменьшается с 92,67 до 88,67\%. Разность -4,00 пункта имеет 95\%-й кластерный интервал [-17,65; +5,83]. Нижняя граница пересекает предел -5, поэтому неухудшение качества не подтверждено. Наличие нуля в интервале также не позволяет утверждать доказанное статистически значимое ухудшение. Отсутствие доказательства различия не является эквивалентностью.


Средняя задержка уменьшается на 1,838 секунды (68,55\%), p95 --- на 1,779 секунды (53,27\%). Парный интервал разности средней задержки [-3,660; -0,079] секунды. Таймер включает маршрутизацию, инструменты и HTTP; холодные загрузки 14,778 и 8,767 секунды и разрыв между частями 998,585 секунды исключены. Эти условия ограничивают перенос результата на другое обслуживание и оборудование.


Всего зарегистрированы 310 локальных и 270 облачных попыток. Облачные ошибки: 12 read timeouts, 6 ошибок JSON, 2 ValueError и один незавершённый ответ. При разном количестве обращений сравнение надёжности политик не доказывает большей надёжности отдельной модели. Пиковая суммарная GPU-память --- 13 758 MiB, RAM системы --- 24,754 GiB по 2862 отсчётам. Это показатели всего устройства и системы с другими приложениями, а не приватная память процесса.


Известная тарифная оценка облачной части hybrid --- \$0,00079139, TOOLS+CLOUD --- \$0,00433758 плюс четыре случая неизвестного usage. Для CLOUD известны \$0,00721649 плюс восемь неизвестных случаев. Эти суммы вычислены из тарифов и returned usage; они не являются счётом об оплате. Электричество, амортизация и полная стоимость локальной эксплуатации не измерены. Неизвестные расходы не заменены нулём.


\section*{6. Диагностика и дополнительные DEV-результаты}

Локальная модель получила 0/30 в численных задачах. Версионный аудит 3 октября независимо пересчитал эталоны с Decimal: 30/30 совпадают с заданными формулами, тогда как все 30 численных ответов выходят за допуск. Ошибок единиц, action, схемы, парсинга, незавершённости или оценщика в этой группе не обнаружено. Промежуточные вычисления отсутствуют; нельзя установить для каждого случая, ошиблась ли модель в арифметике или выборе формулы. Исторические оценки не пересчитывались новым методом.


В compound-задачах hybrid и offline приняли по 12 неверных ветвей. Слова ALERT и OK присутствовали в справочном контексте, и буквальная проверка считала ошибочную метку подтверждённой. Это механизм ложного принятия, а не ошибка вычисленного gold. Новая DEV-версия проверяет формальные условия, отрицание, сущность, единицу и полноту компонентов; к историческим и рабочим политикам она не подключена. Для произвольного текста без проверяемого контракта она возвращает UNKNOWN; экспертизу агронома не заменяет. Типизированные факты должны извлекаться из публичного источника независимо от скрытого gold; объяснения ASK/ABSTAIN требуют согласованного эталона или ручной смысловой проверки.


В отдельном наблюдавшемся DEV36 rule-first получил 31/36, Jev-first --- 30/36, Laya-first --- 23/36. Jev уменьшил генеративные облачные вызовы с 23 до 20, но увеличил все внешние запросы с 23 до 56 за счёт hosted routing. Его превосходство над простым правилом не установлено. Восьми облачным таймаутам соответствуют разные политики, поэтому меньшую задержку Jev нельзя целиком приписать маршрутизатору. Дополнительный lexical guard одинаков для новых pre-routing политик, но отличается от исторического каскада; эффект guard не выдаётся за эффект Jev.


После частичного изменения настроек OpenRouter выполнен отдельный DEV9: Qwen3.8 27B --- 5/9, Ling3.0 Flash Sante --- 4/9, Dots3-Note --- 4/9. Все 27 ответов имеют HTTP200, корректную JSON-схему и reported cost \$0. Median end-to-end latency: 0,673, 1,163 и 1,763 секунды соответственно. Это прежний строгий оценщик и ранее наблюдавшийся DEV; малый результат не является независимым доказательством общего ранжирования. Mercury в историческом эксперименте не заменён этими моделями.


\section*{7. Ограничения и следующий этап}

Доказательная база ограничена одним компьютером, одной локальной квантизацией, конкретным облачным сервисом и выборкой с преобладанием синтетики. Exact span может отклонять эквивалентные формулировки, а простой контекстный verifier --- принимать неполные или неверные составные ответы. Русский и английский поднаборы связаны переводами; наблюдаемая языковая разность не устанавливает качество открытого консультирования на русском. Переиспользование HOME150 после исправления verifier не станет новым подтверждением.


В старой выборке отсутствует gold ABSTAIN. Offline отказался на 20 задачах выбора инструмента, где gold требовал ANSWER; оценить правильность обоснованного отказа по этим данным невозможно. Следующий набор должен различать ASK\_FOR\_DATA при недостающем измерении и ABSTAIN при отсутствии поддерживающего источника. Не следует добавлять несвязанные примеры safety/refusal ради увеличения доли отказов.


Официальный AgXQA test содержит пересечения документов и контекстов с train/validation и прежними данными проекта; после строгого исключения остаётся одна строка. Review40 представляет 20 RU/EN пар по одному блоку ERA5, созданных для независимой ручной проверки. Ни один из авторов ещё не завершил формы, поэтому human\_review=false. Предметные факты вне компетенции студентов направляются эксперту. Машинная разметка или проверка агентом не подменяет такую экспертизу.


Новый confirmatory protocol находится в состоянии DATA\_BLOCKED: независимые группы, человеческий gold и окончательное согласование методов отсутствуют. Новое подтверждающее исследование не запускалось. Перед его запуском фиксируются источники и права, группы документов/погодных блоков, переводы, промпты, модели, retry policy, scoring, учёт задержки и всех external requests. Настройка по ответам final должна быть запрещена. Код и тесты лицензированы отдельно; лицензия статьи и права сторонних материалов остаются самостоятельными решениями.


\section*{8. Выводы}

В контролируемой конфигурации инструменты обеспечили численные ответы и корректный запрос недостающих входов. Добавление локальной модели существенно уменьшило генеративные облачные обращения и наблюдаемую задержку, но не выполнило заранее установленный критерий неухудшения качества. Таким образом, основной положительный эффект относится к использованию локального каскада для экономии обращений, а сохранение качества пока не доказано.


Главное методическое ограничение --- ложное принятие ошибочных составных ответов буквальным verifier. Исправление проверяется только на development-материале и требует нового независимого исследования. Jev-first также не показал превосходства над rule-first и увеличил общее внешнее взаимодействие. Отрицательные и неопределённые результаты сохраняются как часть воспроизводимого эмпирического исследования.


Генеративный ИИ использовался для подготовки кода, организации источников, расчётной поддержки, чернового текста и перевода. За проверку утверждений, ссылок, прав, эталонов и финального текста отвечают человеческие авторы. Человеческая научная проверка и согласие на подачу ещё не завершены. ИИ не указан как автор; благодарность преподавателю и статус научного руководителя не добавлялись автоматически.


\section*{Литература}\begingroup\fontsize{12}{14}\selectfont\begin{enumerate}
\item Chen, L., Zaharia, M., and Zou, J. \emph{FrugalGPT: How to Use Large Language Models While Reducing Cost and Improving Performance}. arXiv:2305.05176, 2023. \url{https://arxiv.org/abs/2305.05176}
\item Aggarwal, P., et al. \emph{AutoMix: Automatically Mixing Language Models}. arXiv:2310.12963, 2023; version 5, 2025. \url{https://arxiv.org/abs/2310.12963}
\item Ding, D., et al. \emph{Hybrid LLM: Cost-Efficient and Quality-Aware Query Routing}. arXiv:2404.14618, 2024. \url{https://arxiv.org/abs/2404.14618}
\item Ong, I., et al. \emph{RouteLLM: Learning to Route LLMs with Preference Data}. ICLR 2025; arXiv:2406.18665. \url{https://arxiv.org/abs/2406.18665}
\item Dekoninck, J., Baader, M., and Vechev, M. \emph{A Unified Approach to Routing and Cascading for LLMs}. ICML 2025, PMLR 267:12987--13010; arXiv:2410.10347. \url{https://proceedings.mlr.press/v267/dekoninck25a.html}
\item Hu, Q. J., et al. \emph{RouterBench: A Benchmark for Multi-LLM Routing System}. arXiv:2403.12031, 2024. \url{https://arxiv.org/abs/2403.12031}
\item Kossen, J., et al. \emph{Semantic Entropy Probes: Robust and Cheap Hallucination Detection in LLMs}. arXiv:2406.15927, 2024. \url{https://arxiv.org/abs/2406.15927}
\item Wen, B., et al. \emph{Know Your Limits: A Survey of Abstention in Large Language Models}. arXiv:2407.18418, 2024, updated 2025. \url{https://arxiv.org/abs/2407.18418}
\item Yang, X., et al. \emph{PLLaMa: An Open-source Large Language Model for Plant Science}. arXiv:2401.01600, 2024. \url{https://arxiv.org/abs/2401.01600}
\item Yang, B., et al. \emph{AgriGPT: a Large Language Model Ecosystem for Agriculture}. arXiv:2508.08632, 2025. \url{https://arxiv.org/abs/2508.08632}
\item Samuel, D. J., et al. \emph{AgroLLM: Connecting Farmers and Agricultural Practices through Large Language Models for Enhanced Knowledge Transfer and Practical Application}. arXiv:2503.04788, 2025. \url{https://arxiv.org/abs/2503.04788}
\item Yang, B., et al. \emph{AgriAgent: Contract-Driven Planning and Capability-Aware Tool Orchestration in Real-World Agriculture}. arXiv:2601.08308, 2026. \url{https://arxiv.org/abs/2601.08308}
\item Ye, Z., et al. \emph{AgroTools: A Benchmark for Tool-Augmented Multimodal Agents in Agriculture}. arXiv:2605.22366, 2026. \url{https://arxiv.org/abs/2605.22366}
\item Yao, S., et al. \emph{tau-bench: A Benchmark for Tool-Agent-User Interaction in Real-World Domains}. arXiv:2406.12045, 2024. \url{https://arxiv.org/abs/2406.12045}
\item Liu, W., et al. \emph{ToolACE: Winning the Points of LLM Function Calling}. ICLR 2025; arXiv:2409.00920. \url{https://arxiv.org/abs/2409.00920}
\item Ma, S., et al. \emph{The Era of 1-bit LLMs: All Large Language Models are in 1.58 Bits}. arXiv:2402.17764, 2024. \url{https://arxiv.org/abs/2402.17764}
\item Gemma Team, et al. \emph{Gemma 4 Technical Report}. arXiv:2607.02770, 2026. \url{https://arxiv.org/abs/2607.02770}
\item Kpodo, J., Kordjamshidi, P., and Nejadhashemi, A. P. \emph{AgXQA: A benchmark for advanced Agricultural Extension question answering}. Computers and Electronics in Agriculture 225:109349, 2024. \url{https://doi.org/10.1016/j.compag.2024.109349}
\end{enumerate}\endgroup
\end{document}
